% Chapter 5 -- Conclusion and Future Work (O-RAN / BMPP-DQN track)
%
% Governs: manuscript/ORAN_BMPP_DQN_Concept_Note_v1.md (Sections 4.2, 6.10,
% 6.11, 6.12, 7, 10's needs-validation flags)
% Chapter mapping: docs/oran_thesis_guide.md
% Status: DRAFT. Intended to be \input{} from the thesis-wide main.tex once
% that skeleton exists; a thesis-wide references.bib does not yet exist, so
% citations below are numbered locally (IEEE style) against the short
% reference list at the end of this file. Renumber/merge into the
% thesis-wide bibliography once main.tex is created.
%
% This chapter closes the \ref{chap:oran-conclusion} forward-references
% made throughout Chapters 1, 3, and 4. Per this track's own academic-
% integrity convention (Chapter 1's own note after its Research Objectives),
% Objective 4's original >=15% framing is restated here exactly as it was
% investigated, not retroactively reworded to match the evidence gathered
% after the fact (HARKing); this chapter states plainly where it was and
% was not met, per Concept Note Section 6.10's own explicit reconciliation.

\chapter{Conclusion and Future Work}
\label{chap:oran-conclusion}

\section{Summary of Contributions}
\label{sec:oran-summary-contributions}

This thesis designed, implemented, and evaluated BMPP-DQN (Branching
Multi-Pass Parameterized DQN), a deep reinforcement learning architecture
combining branching discrete-action decomposition~\cite{tavakoli2018branching}
with multi-pass parameterized continuous-action
processing~\cite{xiong2018pdqn,bester2019mpdqn} within a single coupled
network, for joint RU activation, functional split selection, transmit
power control, and PRB allocation in a simulated O-RAN energy-optimization
scenario. Four contributions, each evidenced in
Chapters~\ref{chap:oran-system-model}--\ref{chap:oran-results}, follow from
this design:

\begin{enumerate}
  \item \textbf{An architectural integration not previously proposed}
    (Section~\ref{sec:oran-significance}): BMPP-DQN couples branching's
    $R$-independent discrete decomposition with multi-pass processing's
    principled discrete-continuous parameter coupling, closing the gap
    identified in Section~\ref{sec:oran-problem-statement} between OREO's
    flat, thresholded continuous relaxation~\cite{qazzaz2026oreo} and the
    two architectures this design draws on, neither of which alone handles
    O-RAN's hybrid, multi-timescale action space.
  \item \textbf{A formal parameterized MDP formulation}
    (Section~\ref{sec:oran-mdp}) spanning four action branches across two
    decision timescales, with a fronthaul- and switching-aware reward
    function (Eq.~\ref{eq:oran-reward}) and a literature-informed power
    model (Section~\ref{sec:oran-power}).
  \item \textbf{A four-method comparison under matched conditions}
    (Chapter~\ref{chap:oran-results}) against DQN, DDPG, and MP-DQN,
    cross-checked by three independent multi-criteria decision-analysis
    schemes (equal-weighted TOPSIS, entropy-weighted TOPSIS, and VIKOR;
    Sections~\ref{sec:oran-mcda}--\ref{sec:oran-mcda-robustness}), a
    power-model sensitivity sweep across 9 configurations
    (Section~\ref{sec:oran-sensitivity}), and a supplementary $n=10$-seed
    statistical-power validation (Section~\ref{sec:oran-n10}) -- a more
    extensive robustness-checking protocol than this thesis's own original
    scope required, undertaken specifically to avoid over-claiming from a
    single table of $n=3$ results.
  \item \textbf{A direct empirical answer to RQ3}
    (Section~\ref{sec:oran-rq3}): the two-timescale design's measurable
    contribution is discrete-decision stability (a $6\times$ reduction in
    switching frequency), not raw energy or QoS performance, resolving a
    research question this track's concept note had posed but left
    empirically untested until this thesis's own ablation.
\end{enumerate}

\section{Revisiting the Research Objectives}
\label{sec:oran-revisit-objectives}

Section~\ref{sec:oran-objectives} stated five research objectives. Objectives
1-3 (design BMPP-DQN; formulate the parameterized MDP; implement the
focused single-gNB environment) are met in full, as documented in
Chapter~\ref{chap:oran-system-model}. Objective 5 (analyze the
two-timescale design's effect on convergence and energy efficiency) is
answered directly in Section~\ref{sec:oran-rq3}: the design improves
decision stability without a measurable convergence or efficiency cost in
either direction.

Objective 4 -- quantify BMPP-DQN's energy-saving performance against DQN,
DDPG, and MP-DQN, against a target of $\geq15\%$ savings relative to these
baselines -- is stated here, as Section~\ref{sec:oran-objectives} already
flagged it would be, without qualification: \textbf{this target is not met
by the evidence gathered}, at $n=3$ or at the supplementary $n=10$
validation. BMPP-DQN's mean power is not $15\%$ below any baseline's; it is
essentially tied with MP-DQN's, the lowest-power baseline, at both sample
sizes (Table~\ref{tab:oran-convergence}, Table~\ref{tab:oran-convergence-n10}).
This objective is restated here exactly as it was originally planned and
investigated, rather than retroactively reworded to match the evidence
gathered after the fact, per this track's own stated discipline
(Section~\ref{sec:oran-objectives}'s own forward-reference to this
chapter).

What the evidence \emph{does} support, in place of a raw energy-savings
margin, is detailed in Section~\ref{sec:oran-rq-answers}: BMPP-DQN is the
strongest or joint-strongest method under every multi-criteria scheme
tested, at both sample sizes, a finding independently robust to an
order-of-magnitude perturbation of every unvalidated power-model constant
(Section~\ref{sec:oran-sensitivity}). This is a genuine, evidenced
contribution -- a multi-criteria balance across power, both QoS metrics,
and switching stability simultaneously -- but it is a different claim in
kind from a demonstrated energy-savings percentage, and the two are not
conflated here or in Section~\ref{sec:oran-significance-revisited} below.

\section{Significance of the Research, Revisited}
\label{sec:oran-significance-revisited}

Section~\ref{sec:oran-significance} stated this thesis's intended
significance before its results were known. Revisited in light of
Chapter~\ref{chap:oran-results}:

\begin{itemize}
  \item \textbf{Architectural contribution} (DRL literature on hybrid
    action spaces): unaffected by the results -- the integration of
    branching and multi-pass processing into a single coupled network is a
    design contribution independent of how it ultimately performs, and
    Objective 1/RQ1 (Section~\ref{sec:oran-rq-answers}) confirm it trains
    stably across every seed tested.
  \item \textbf{O-RAN application relevance}: unaffected -- the problem
    formulation (Chapter~\ref{chap:oran-system-model}) remains a faithful
    structural match to O-RAN's disaggregated RU/DU/CU control surface and
    its Non-RT/Near-RT RIC timescale split, independent of this specific
    algorithm's comparative performance.
  \item \textbf{Energy savings}: \emph{not demonstrated}
    (Section~\ref{sec:oran-revisit-objectives}). This bullet is retained,
    not deleted, because it was this thesis's original stated significance
    claim and the evidence against it is itself part of this thesis's
    record, not something to omit.
  \item \textbf{Multi-criteria balance} (evidenced): BMPP-DQN is the most
    balanced of the four methods across power, strict QoS, per-UE QoS,
    switching stability, and throughput simultaneously, a finding that
    survives a power-model sensitivity sweep and that -- with the
    qualification in Section~\ref{sec:oran-n10} that its strongest form
    (full three-scheme agreement) is partly an $n=3$-specific result -- also
    survives the move to $n=10$ seeds under at least one of the three MCDA
    schemes tested in every case. This, not the energy-savings bullet
    above, is this thesis's actual, defensible empirical contribution
    relative to RQ2.
\end{itemize}

\section{Limitations}
\label{sec:oran-limitations}

Section~\ref{sec:oran-scope} bounded this thesis's scope in advance; those
same bounds are this chapter's limitations, read in light of what was
actually found:

\begin{itemize}
  \item \textbf{Simulation-only evaluation.} All results use a simplified
    log-distance-path-loss-plus-Rayleigh-fading channel model
    (Section~\ref{sec:oran-channel}) and a literature-informed but not
    independently validated power-consumption model
    (Section~\ref{sec:oran-power-validation}). No real O-RAN testbed,
    hardware power measurement, or over-the-air channel was used. The
    power-model sensitivity sweep (Section~\ref{sec:oran-sensitivity})
    shows the \emph{comparative} findings are insensitive to this model's
    unvalidated constants' order of magnitude, but this does not establish
    that the \emph{absolute} power, reward, or QoS figures reported in this
    chapter would transfer to a physical deployment.
  \item \textbf{Single-gNB, downlink-only scope.} The environment models one
    gNB with $R=4$ RUs, one DU, and one CU (Section~\ref{sec:oran-topology});
    no multi-cell coordination, inter-cell interference beyond the
    single-pool model already included, or uplink traffic is considered.
  \item \textbf{A tractable, 3-option split abstraction.} The three
    centralization levels used here approximate 3GPP TR~38.801 Options 2, 6,
    and 8 (Section~\ref{sec:oran-topology}) rather than the full space of
    standardized functional splits.
  \item \textbf{Unvalidated reward weights.} The reward function's three
    weights ($\alpha$, $\beta_{\mathrm{qos}}$, $\gamma_{\mathrm{sw}}$,
    Section~\ref{sec:oran-mdp}) were chosen to preserve a reasonable
    relative magnitude, not derived from a literature source or tested in a
    sensitivity sweep of their own -- unlike the power-model constants, which
    were swept in Section~\ref{sec:oran-sensitivity}. Given that BMPP-DQN's
    per-UE QoS deficit relative to the baselines was investigated across
    several training-mechanics explanations without this one being tested
    (the independent-exploration fix, PER, and curriculum learning were all
    ruled in or out; reward-weight sensitivity was not), this remains the
    most likely untested lever on the energy/QoS/switching trade-off
    reported in Chapter~\ref{chap:oran-results}.
  \item \textbf{Thin statistical power at the canonical scale.} The
    supervisor-approved protocol's $n=3$ seeds is a thin basis for the
    paired significance tests reported throughout
    Chapter~\ref{chap:oran-results}; the supplementary $n=10$ validation
    (Section~\ref{sec:oran-n10}) was undertaken specifically to probe this,
    and it changed which comparisons reach significance (DDPG's reward
    deficit) and whether the headline multi-criteria ranking holds under
    every scheme simultaneously (it does not, at $n=10$, under
    equal-weighted TOPSIS specifically). Both the $n=3$ and $n=10$ results
    are reported in this thesis rather than only the more favorable of the
    two.
  \item \textbf{Three baselines, not an exhaustive comparison.} BMPP-DQN is
    compared against DQN, DDPG, and MP-DQN
    (Section~\ref{sec:oran-scope}); OREO~\cite{qazzaz2026oreo} is discussed
    only qualitatively (Section~\ref{sec:oran-problem-statement}), and
    HyAR~\cite{li2022hyar}, a VAE-based hybrid-action alternative considered
    but not adopted for this architecture, is likewise not reproduced as a
    baseline.
\end{itemize}

\section{Future Work}
\label{sec:oran-future-work}

Four directions follow directly from the limitations above and from
findings in Chapter~\ref{chap:oran-results} that this thesis's own scope
could investigate but not resolve.

\textbf{Validate the power model and reward weights against real
measurements.} Section~\ref{sec:oran-power-validation}'s needs-validation
constants (the absolute RU/DU/CU/fronthaul wattage figures) and
Section~\ref{sec:oran-mdp}'s reward weights are the two remaining
unvalidated numerical components of this thesis's model. The power-model
sensitivity sweep (Section~\ref{sec:oran-sensitivity}) establishes that the
comparative ranking is insensitive to the former's order of magnitude;
whether the same holds for the latter, and whether either can be replaced
with commercial O-RAN hardware measurements, is open. A real-testbed study
in the style of Bordin et al.'s ns-O-RAN-based evaluation~\cite{bordin2025design}
would address both simultaneously.

\textbf{Resolve the ranking volatility surfaced at $n=10$.}
Section~\ref{sec:oran-n10} found that BMPP-DQN's full three-scheme MCDA
agreement, this thesis's strongest robustness claim, partly depends on
sample size. A larger validation run (beyond this thesis's $n=10$,
approaching the kind of sample size a dedicated follow-up study could
afford) would establish whether DDPG's equal-weighted-TOPSIS lead at
$n=10$ is itself an artifact of a still-limited sample, or a genuine,
larger-sample-confirmed exception to BMPP-DQN's otherwise-consistent
multi-criteria strength.

\textbf{Extend beyond a single gNB.} This thesis's single-agent,
single-gNB scope was a deliberate tractability bound
(Section~\ref{sec:oran-scope}), justified in part by Liang et al.'s
federated-TD3 O-RAN energy-optimization work~\cite{liang2026federated}
showing a feasible path to multi-cell coordination via an rApp
aggregator plus per-cell xApp agents, rather than a single centralized
policy. Extending BMPP-DQN's branching-plus-multi-pass design to that
federated setting -- one BMPP-DQN instance per cell, aggregated rather than
jointly trained -- is a natural next step that preserves this thesis's core
architectural contribution while removing its single-gNB limitation.

\textbf{Test functional-split abstractions beyond the 3-option mapping.}
Section~\ref{sec:oran-topology}'s three-level split abstraction is a
tractability simplification whose bandwidth ratios do not match 3GPP
TR~38.801's real per-option figures exactly
(Section~\ref{sec:oran-power-validation}'s own disclosure). A follow-up
could use the full 8-option split space directly, at the cost of a larger
discrete action space than this thesis's branching architecture was sized
for, testing whether BMPP-DQN's branching decomposition scales to that
larger space without the flat baselines' (MP-DQN's, in particular)
combinatorial tractability ceiling.

\section{Closing Statement}
\label{sec:oran-closing}

This thesis set out to answer whether branching action decomposition and
multi-pass parameterized processing could be combined into a single DRL
architecture for O-RAN energy optimization, and whether doing so would
yield energy savings over existing discrete, continuous, and parameterized
baselines. The first question is answered affirmatively: BMPP-DQN trains
stably, end-to-end, as the proposed coupled architecture
(Section~\ref{sec:oran-summary-contributions}). The second is answered in
the negative, stated plainly rather than reframed after the fact
(Section~\ref{sec:oran-revisit-objectives}). What the evidence gathered in
its place supports -- a multi-criteria balance across power, QoS, and
switching stability, confirmed robust to this model's own unvalidated
constants and substantially, though not completely, robust to a larger
statistical sample -- is a narrower contribution than originally targeted,
but one this thesis can defend with the evidence actually in hand, which is
the standard Chapter~\ref{chap:oran-results} and this chapter have tried to
hold themselves to throughout.

% ---------------------------------------------------------------------------
% This chapter's citations (qazzaz2026oreo, xiong2018pdqn, bester2019mpdqn,
% tavakoli2018branching, bordin2025design, liang2026federated, li2022hyar)
% resolve against the consolidated thesis-wide bibliography in
% thesis/chapters/references_oran.tex, not a local thebibliography here.
% ---------------------------------------------------------------------------
